% Cover letter based on main.tex. Compile with pdfLaTeX.
\documentclass[11pt,a4paper]{article}
\usepackage[top=25mm,bottom=25mm,left=25.4mm,right=25.4mm]{geometry}
\usepackage[T1]{fontenc}
\usepackage{newtxtext,newtxmath}
\usepackage{microtype}
\usepackage[hidelinks]{hyperref}
\pagestyle{empty}
\setlength{\emergencystretch}{1em}
\setlength{\parindent}{14.5pt}
\setlength{\parskip}{9pt}
\AtBeginDocument{\fontsize{11}{14.5}\selectfont}

\newcommand{\LetterDate}{\today}
\newcommand{\JournalName}{Journal of Applied Crystallography}
\newcommand{\ManuscriptTitle}{Polarized neutron diffraction with \textit{ex-situ} $^3$He neutron spin filter and two-dimensional detector for a complex incommensurate magnetic order}
\newcommand{\AuthorName}{Shingo Takahashi}
\newcommand{\Affiliation}{Institute for Solid State Physics, The University of Tokyo}
\newcommand{\PostalAddress}{5-1-5 Kashiwanoha, Kashiwa, Chiba 277-8581, Japan}
\newcommand{\PhoneNumber}{+81-29-287-8910}
\newcommand{\EmailAddress}{shingo.takahashi@issp.u-tokyo.ac.jp}

\begin{document}
{\raggedleft \LetterDate\par}
\vspace{5pt}
\noindent Dear Editors of \JournalName,

% 第1段落：研究の紹介、背景、従来手法、残された課題。
The present manuscript reports an experimental study on polarized neutron diffraction using a $^3$He neutron spin filter and a two-dimensional detector to investigate complex incommensurate magnetic orders.
Polarized neutron scattering has long been used to determine magnetic moment directions and distinguish magnetic scattering from nuclear scattering.
Recent interest in multiferroics, magnetic skyrmions, and related magnetic orders has increased the need to probe magnetic satellite reflections distributed in three-dimensional reciprocal space.
Conventional triple-axis setups with Heusler-crystal analyzers, however, are generally restricted to reflections within the horizontal scattering plane.
A $^3$He neutron spin filter can analyze neutrons scattered over a finite solid angle and is therefore well suited to measurements with a two-dimensional detector.
This combination has been used in small-angle neutron scattering, although the accessible momentum-transfer range is generally limited to low $Q$.
Extending this approach beyond the small-angle regime requires placing the spin filter and detector at finite scattering angles.
Related instrumental concepts have been reported, but quantitative experimental demonstrations of polarization analysis with such setups remain limited.

% 第2段落：研究の焦点、装置と補正法、検証、物質への適用と成果。
In the present study, we focus on quantitative polarization analysis of magnetic reflections outside the horizontal scattering plane.
We implemented an experimental setup combining an \textit{ex-situ} $^3$He neutron spin filter and a two-dimensional detector at the 5G beamport of JRR-3.
We established a procedure to reconstruct three-dimensional reciprocal-space intensity distributions while correcting each detector pixel for imperfect neutron polarization and time-dependent analyzer transmission.
The correction parameters were determined from a nuclear Bragg reflection and NMR monitoring, and the correction was validated by recovering the expected non-spin-flip character of nuclear scattering.
We then applied this method to single-crystal diffraction measurements on the multiferroic Mn$_3$WO$_6$ at 23~K, whose three symmetry-related magnetic propagation vectors cannot all lie in a single plane through the reciprocal-space origin.
The corrected intensities revealed predominantly non-spin-flip scattering at one incommensurate reflection and spin-flip scattering at another.
Their measured asymmetries are inconsistent with moments parallel to the $c$ axis and are reproduced by an oblique sinusoidal model, demonstrating the sensitivity of the method to magnetic moment directions.

% 第3段落：学術的意義と投稿誌への適合性。
The present results provide an experimentally validated approach to quantitative polarization analysis in three-dimensional reciprocal space.
These advances in polarized neutron diffraction are expected to investigate complex incommensurate magnetic structures.
We believe that the present manuscript meets the criteria for publication in \textit{\JournalName}.

\vspace{4pt}
\noindent\begin{minipage}{\textwidth}
\raggedleft
Yours sincerely,\\[4pt]
\AuthorName\\[7pt]
\Affiliation\\
\PostalAddress\\
Tel: \PhoneNumber\\
\href{mailto:\EmailAddress}{\EmailAddress}
\end{minipage}
\end{document}
